% !TEX root = tarasov-qasp-zh.tex
% 本章英文源（ch-14.tex）内无 sec 扉图块与章首格言（有效代码），故中文版照原样不设。
\chapter[关于物体的平衡，你知道些什么？]{关于物体的平衡，你知道些什么？\label{ch-14}}

\begin{dialogue}

\tchr \figr{fig-59} 画出了砖的两个平衡位置。两个位置都是稳定平衡（stable equilibrium），但稳定的程度不同。\hlred{这两个位置中哪一个更稳定？}

\begin{figure}[!h]%
\centering
\includegraphics[width=0.4\textwidth]{figs/ch-14/fig-059a.pdf}
\caption[哪块砖更稳定？]{哪块砖更稳定？\label{fig-59}}
\end{figure}

\stda 显然，是\figr{fig-59}~\drkgry{(a)} 所示的位置。

\tchr 为什么？

\stda 在这个位置，砖的重心离地面更近，

\tchr 这还不是全部。

\stdb 这个位置的支承面（bearing surface）面积比\figr{fig-59}~\drkgry{(b)} 所示位置的大。

\tchr 这同样不是全部。为了把事情弄清楚，我们来考察两个物体的平衡：底面为正方形的长方体和圆柱体（\figr{fig-60}~\drkgry{(a)}）。设长方体和圆柱体的高度同为 $H$，底面积同为 $S$。这样，两个物体的重心位于同一高度，支承面的面积也彼此相同。然而，它们的稳定程度却不相同。

某一具体平衡状态的稳定程度，用为使物体永久地偏离该状态所必须耗费的能量来量度。
\pnote{把重心抬到哪里才算“永久地”扰动？答案是抬到支承面边缘（翻转棱）的正上方：那是翻倒全过程中重心最高的一处，一旦越过，重力矩自会完成剩下的翻倒，无须再做功。因此所需的最小能量等于 $mg$ 与这一抬高量之积，下文的 $h_{1}$、$h_{2}$ 正是这样的抬高量。大学普物对稳定平衡的普遍刻画——势能取极小值，任何扰动都使势能升高——正是这里能量判据的一般形式。}

\begin{figure}
\centering
\includegraphics[width=0.5\textwidth]{figs/ch-14/fig-060a.pdf}
\caption[比较两个物体的平衡。]{比较两个物体的平衡，哪个更稳定？\label{fig-60}}
\end{figure}
\stdb 您用的“永久地”一词是什么意思？

\tchr 意思是：此后若任凭物体自行发展，它再也不能回到初始状态。这份能量等于物体的重力与一个高度的乘积——为使物体回不到初始位置，重心必须被抬升到这个高度。在长方体与圆柱体的例子中，圆柱体的半径为 $R=\sqrt{S/\pi}$，长方体底面的边长为 $a=\sqrt{S}$。要破坏圆柱体的平衡，须把它的重心抬高（\figr{fig-60}~\drkgry{(b)}）
\begin{equation*}%
h_{1} = \sqrt{ \left( \frac{H}{2}\right)^{2} + R^{2}} - \frac{H}{2}
\end{equation*}
要破坏长方体的平衡，须把它的重心抬高（\figr{fig-60}~\drkgry{(b)}）
\begin{equation*}%
h_{2} = \sqrt{ \left( \frac{H}{2}\right)^{2} +  \left(\frac{a}{2}\right)^{2}} - \frac{H}{2}
\end{equation*}
由于
\begin{equation*}%
 \left( \frac{a/2}{R} \right)  = \frac{\sqrt{\pi S}}{2\sqrt{S}} = \frac{\sqrt{\pi}}{2} < 1
\end{equation*}
可得 $h_{2} < h_{1}$。因此，在所考察的两个物体中，圆柱体更稳定。
\pnote{同底面积的圆柱为何更稳？物体总沿“最容易”的方向翻倒，即绕使重心抬高量最小的那条支承面边缘翻转。圆柱底面各方向的支撑半径都等于 $R$；方形底面沿垂直于边的方向只有 $a/2$（沿对角线方向虽大，不起作用）。由 $a/2=\left(\sqrt{\pi}/2\right)R\approx 0.89R$ 知，长方体的“最短支撑半径”更小，重心要翻越的门槛更低，故更容易被翻倒。重心同高、底面积相同而稳定性不同——这正说明几何比较代替不了能量比较。}

现在我们回到砖的两个位置那个例子上来。

\stda 若把砖一路翻过去，它将依次从一个平衡位置进入另一个平衡位置。\figr{fig-61} 中的虚线画出了这个过程中砖的重心描出的轨迹。要改变砖的平放位置，须把它的重心抬高 $h_{1}$，耗费的能量等于 $mgh_{1}$；要改变它的竖放位置，须把重心抬高 $h_{2}$，耗费的能量为 $mgh_{2}$。平放的砖稳定程度更大，原因在于
\begin{equation}%
mgh_{1} > mgh_{2}
\label{eq-82}
\end{equation}

\begin{figure}[!h]%
\centering
\includegraphics[width=0.4\textwidth]{figs/ch-14/fig-061a.pdf}
\caption[砖翻倒过程中其重心描出的轨迹。]{砖翻倒过程中其重心描出的轨迹。\label{fig-61}}
\end{figure}
\tchr 你总算把物体平放时更稳固这一点论证出来了。

\stdb 可是，$h_{1}$ 和 $h_{2}$ 显然取决于重心离地面的高度和底面积的大小。这是不是意味着，讨论物体的稳定程度时，比较重心的高低和底面积的大小就是正确的做法呢？

\tchr 是啊，可以这么说，但只在这些量影响到高度 $h_{1}$ 与 $h_{2}$ 之差的范围内成立。比如在长方体与圆柱体的例子中，比较重心的高低和底面积的大小，并不足以判定哪个物体更稳定。此外，我想请你们注意下面一点。到目前为止，我们一直默认两个物体由同种材料制成。在这种情形，只要满足几何条件 $h_{1}>h_{2}$，不等式~\eqref{eq-82} 便成立。

然而在一般情形下，物体可以由不同材料制成；由于各物体密度不同，甚至当 $h_{1}<h_{2}$ 时，不等式~\eqref{eq-82} 也可能得到满足。例如，平放的软木砖就不如竖放的铅砖稳定。现在来看看：你们知道物体的平衡要满足哪些条件？
\pnote{这段话的关键在比较对象。软木密度约 $0.2\,\mathrm{g/cm^{3}}$，铅约 $11.3\,\mathrm{g/cm^{3}}$，同尺寸的铅砖质量是软木砖的五十余倍。于是平放的软木砖尽管 $h_{1}$ 较大，$mgh_{1}$ 仍远小于竖放铅砖的 $mgh_{2}$——这是不等式~\eqref{eq-82} 因密度悬殊而不成立的例子。而上一句“即使 $h_{1}<h_{2}$，式~\eqref{eq-82} 也能满足”说的是相反的搭配：密实的物体平放、轻小的物体竖放。两种情形是同一个教训：材料不同时，比较稳定程度必须直接比较能量，不能停留在几何量上。}

\stda 作用在物体上的所有力之和应为零。此外，物体的重力矢量应落在其底面的范围之内。
\pnote{学生甲的第一条即共点力的平衡条件：合力为零（人教版必修第一册第三章“共点力的平衡”一节）。第二条的确切含义是重力作用线不越出支承面：一旦越出，重力对支承面边缘的力矩就会把物体翻倒——这正是本章开头砖块问题的判据。}

\tchr 好。不过，更好的做法是把平衡条件表述成另一种更普遍、更便于实际应用的形式。应当区分两条平衡条件：
\begin{description}[leftmargin=1cm,font=\bfseries]
\item[第一个条件：] 作用在物体上的所有力沿任一方向的投影应当彼此抵消。换句话说，所有各力沿任一方向的投影的代数和应为零。利用这条条件可以写出的方程数，等于问题中独立方向的数目：一维问题一个方程，二维问题两个方程，一般情形三个方程（所取方向相互垂直）。

\item[第二个条件（力矩条件）：] 各力对任一点的力矩（moment of a force）的代数和应为零。其中，凡使物体绕所选点沿某一方向（譬如顺时针方向）转动的力矩取正号，凡使物体沿相反方向（逆时针方向）转动的力矩取负号。具体列出力矩条件时，按以下步骤进行：
\begin{enumerate}[label=(\alph*), leftmargin=1cm]
 \item 确定作用在物体上的全部力；
 \item 选取一个点，各力的力矩将对这一点来计算；
\item 求出所有各力对所选点的力矩；
 \item 写出力矩代数和的方程，并令它等于零。
 \end{enumerate}
\end{description}
 应用力矩条件时，应当记住：
\begin{enumerate}[label=(\arabic*), leftmargin=1cm]
\item 上面所说的条件，属于问题中所有各力及其力臂（arm）都处于同一平面内的情形（问题不是三维的）；
\item 力矩的代数和应对任一点令其为零，这个点既可取在物体内部，也可取在物体外部。
\end{enumerate}
\pnote{力矩、力臂和力矩平衡条件已不在现行高中人教版教材的要求之内；中学最接近的内容是初中“简单机械”一章的杠杆平衡条件（动力×动力臂＝阻力×阻力臂，人教版八年级下册），它相当于只含两个力的力矩方程。本节的两条平衡条件合在一起，就是大学普通物理中的平面力系平衡方程；梁与支座反力一类问题超出现行高考范围，在物理竞赛中则是常规工具。}

应当强调：虽然各个力矩的数值确实与所选的点（力矩对它来计算）有关，力矩的代数和却无论如何都等于零。为了更好地理解平衡条件，我们来考虑一个具体的问题。

\hlred{一根重力为 $P_{1}$ 的梁（beam）支承在 $B$、$C$ 两点（\figr{fig-62}~\drkgry{(a)}）。在点 $D$ 处，梁上挂有一个重力为 $P_{2}$ 的重物。距离 $\overline{AB}=a$、$\overline{BC}=2a$、$\overline{CD}=a$。求两个支座处的支座反力（reaction）$N_{B}$ 和 $N_{C}$。设支座反力沿竖直方向。}

和通常一样，先指明作用在物体上的力。
 \begin{figure}[!htb]
\centering
\includegraphics[width=.6\textwidth]{figs/ch-14/fig-062a.pdf}
\caption[带悬挂重物的梁。]{带悬挂重物的梁，问题是求支座反力。\label{fig-62}}
\end{figure}
\stda 本题中的物体是梁。它受到四个力：重力 $P_{1}$ 与 $P_{2}$，支座反力 $N_{B}$ 与 $N_{C}$。

\tchr 把这些力画到图上。

\stda 可我不知道反力是朝上还是朝下。

\tchr 先假定两个反力都朝上。

\stda 好，这是我的图（\figr{fig-62}~\drkgry{(b)}）。接着我可以把第一个平衡条件写成方程
\begin{equation*}%
N_{B}+N_{C}=P_{1}+P_{2}
\end{equation*}

\tchr 对这个方程本身我没有异议。不过在本题中，用第二个平衡条件（力矩条件）更简单：先对点 $B$ 用它，再对点 $C$ 用它。

\stda 好吧，我就这样做。结果写出了方程
\begin{equation}
\left.
\begin{aligned}
 \text{对点 $B$：} aP_{1}-2aN_{C}+3aP_{2} &= 0  \label{eq-83}\\
 \text{对点 $C$：} 2aN_{B}-aP_{1}+aP_{2} &=0
 \end{aligned}
 \right\}
\end{equation}

\tchr 现在你看到了：你的每个方程里都只含一个未知量，很容易求出。
\pnote{值得点破教师选矩心的用意：对点 $B$ 取矩时，反力 $N_{B}$ 的力臂为零，根本不出现在第一个方程里；对点 $C$ 取矩时，$N_{C}$ 同样消失。把矩心（力矩的参考点）选在未知力的作用线上，一个方程便只含一个未知量——这是力矩法解这类问题的核心技巧。}
\stda 由方程组~\eqref{eq-83} 解得
\begin{align}%
N_{B} =\frac{\left( P_{1}-P_{2} \right)}{2}
\label{eq-84}\\
N_{C} =\frac{\left( P_{1}+3P_{2} \right)}{2}
\label{eq-85}
\end{align}

\tchr 式~\eqref{eq-85} 恒取正值。这就是说，反力 $N_{C}$ 总是朝上的（与我们所作的假定一致）。式~\eqref{eq-84} 在 $P_{1}>P_{2}$ 时取正值，$P_{1}<P_{2}$ 时取负值，$P_{1}=P_{2}$ 时等于零。这就是说：当 $P_{1}>P_{2}$ 时，反力 $N_{B}$ 与我们假定的方向一致，即朝上（见\figr{fig-62}~\drkgry{(b)}）；当 $P_{1}<P_{2}$ 时，反力 $N_{B}$ 朝下（见\figr{fig-62}~\drkgry{(c)}）；而当 $P_{1}=P_{2}$ 时，反力 $N_{B}$ 不存在。
\pnote{底本式~\eqref{eq-85} 左端误排为 $N_{B}$（与式~\eqref{eq-84} 相同），按教师随后的解读（该式恒正、给出反力 $N_{C}$）应为 $N_{C}$；又，“当 $P_{1}>P_{2}$ 时，反力 $N_{B}$ 与假定的方向一致（朝上）”一句中，底本不等号误排为 $P_{1}<P_{2}$，与前句“式~\eqref{eq-84} 在 $P_{1}>P_{2}$ 时取正值”自相矛盾。两处均按上下文径改；经核对本书所据的 2020 年英文重排印本，两处亦然，是重排之误抑或承自更早版本，待考。另式~\eqref{eq-83} 中 $N_{C}$ 下标误作小写及个别标点级排印小误，均径改不注。}

\end{dialogue}


\cleardoublepage

% NEWTERMS: equilibrium of bodies | 物体的平衡 | 本章主题；equilibrium position（平衡位置）已见于第 10 章
% NEWTERMS: stable equilibrium | 稳定平衡 | 原书仅说 positions are stable，未展开稳定/不稳定/随遇三类
% NEWTERMS: degree of stability | 稳定程度 | 俗称“稳度”；原书无专门短名
% NEWTERMS: bearing surface | 支承面 | 与第 2 章 bearing reaction（支承力）同源构词，注意与“底面（base）”区分
% NEWTERMS: rectangular parallelepiped (with a square base) | 长方体（底面为正方形的） | 普通几何词；第 15 章重心问题恐再现，留意泛称“平行六面体”
% NEWTERMS: right circular cylinder | 圆柱体 | 普通几何词，即正圆柱
% NEWTERMS: moment of a force / force moment | 力矩 | 第 2 章 rotational moment 译“转动矩”（力偶矩语境），本章为标准语境
% NEWTERMS: arm (of a force) | 力臂 | 沿用初中杠杆术语
% NEWTERMS: moment condition | 力矩条件 | 原书第二个平衡条件名称
% NEWTERMS: beam | 梁 | 力学语境首现；现行高中教材不出现
% NEWTERMS: reaction (of a support) | 支座反力 | 区别于第 2 章 bearing reaction（支承力/支持力）
% NEWTERMS: permanently disturb | 永久（地）扰动 | 原书量度稳定程度的关键限定词
